% ---
% id: protocol-compiler/ex:transcript
% kind: ex
% title: "A three-message transcript"
% summary: >-
%   Alice announces the kind of round, Bob discloses his bit only on test
%   rounds, and Alice then discloses hers only on detected test rounds. Party
%   computability holds by the signatures of the tables, and two independent
%   announcement maps cannot express the protocol.
% prominence: standard
% status: accepted
% uses: []
% mentions: [def:declaration, thm:sifting-and-key, transmitter/ass:seq]
% assumes: []
% refs: []
% provenance: newly-derived
% lean:
% notes: "The timing premise it names is the transmitter's A-SEQ."
% ---

\begin{example}\label{ex:transcript}
Default readouts, $\rd{\mathsf A}(x,v)=x$ and $\rd{\mathsf B}(x,v)=v$. Alice's
letter is $x=(z,b)$, with $z\in\{\mathrm{test},\mathrm{gen}\}$ saying whether
the round is a test or a generation round and a bit $b\in\{0,1\}$; Bob's
letter is $v\in\{0,1,\varnothing\}$, a bit or no detection. The transcript
(Definition~\ref{def:declaration}\,(DX5)) has $L=3$ messages:
\begin{enumerate}[label=(m\arabic*),leftmargin=3.0em]
\item $\pi_1=\mathsf A$, $\mathcal{C}_1=\{\mathrm{test},\mathrm{gen}\}$,
  $a_1(z,b)=z$: Alice announces the kind of round;
\item $\pi_2=\mathsf B$,
  $\mathcal{C}_2=\{0,1,\mathrm{click},\varnothing\}$, and
  $a_2(v,\kappa_1)=v$ if $\kappa_1=\mathrm{test}$ or $v=\varnothing$, and
  $a_2(v,\mathrm{gen})=\mathrm{click}$ for $v\in\{0,1\}$: Bob discloses his
  bit on test rounds and only the detection on generation rounds;
\item $\pi_3=\mathsf A$, $\mathcal{C}_3=\{0,1,\ast\}$, and
  $a_3((z,b),\kappa_1,\kappa_2)=b$ if $\kappa_1=\mathrm{test}$ and
  $\kappa_2\ne\varnothing$, and $\ast$ otherwise: Alice discloses her bit on
  detected test rounds.
\end{enumerate}
By \eqref{eq:transcript}, $a((\mathrm{test},b),v)=(\mathrm{test},v,b)$ for
$v\in\{0,1\}$, $a((\mathrm{test},b),\varnothing)
=(\mathrm{test},\varnothing,\ast)$,
$a((\mathrm{gen},b),v)=(\mathrm{gen},\mathrm{click},\ast)$ for $v\in\{0,1\}$
and $a((\mathrm{gen},b),\varnothing)=(\mathrm{gen},\varnothing,\ast)$: seven
classes occur out of the twenty-four of
$\Kann=\mathcal{C}_1\times\mathcal{C}_2\times\mathcal{C}_3$, and the tables
are total on all of their domains as D(viii) asks, their values on
combinations that never occur being immaterial. Sifting passes on
$(\mathrm{gen},\mathrm{click},\ast)$ alone; with $\hold=\mathsf A$ the key
rule is $r^{0}_{\mathsf A}((z,b),\kappa)=b$, a function of Alice's readout.
The acceptance map may send each test class $(\mathrm{test},v,b)$ to the
pair $(v,b)$, so that the test statistics see both bits and the generation
rounds disclose neither.

\emph{What two \tref[announcement map]{announcement maps} cannot express.} In
the \tref{basic format}, with two maps $a_A$ on $\Oset{\mathsf A}$ and $a_B$
on $\Oset{\mathsf B}$, Bob's announcement is a function of $v$ alone and cannot
depend on the kind of round, which he learns only from (m1); the only way to
write this protocol there is to put $z$ into Bob's readout, which misstates
what Bob measures. Here party computability holds \emph{by the signatures}:
$a_2$ has no argument through which $x$ could enter except the message
$\kappa_1$ that Alice has made public (Theorem~\ref{thm:sifting-and-key}\,(iv)).

\emph{What it is read under.} Every message of a round follows every
measurement of that round (\xref{transmitter/ass:seq}): (m1) is sent after Bob
has measured, so his measurement cannot depend on it, and the declaration has
no slot in which such a dependence could be written, $\{F_v\}$ being one fixed
family. A protocol in which Bob chose his measurement after hearing $z$ is
outside this format and outside the fixed-measurement \tref{round model}.
\end{example}
